\begin{figure}[tbp]
\centering
\begin{tikzpicture}[
  box/.style={draw=MidnightBlue,rounded corners=2pt,align=center,
    text width=4.05cm,minimum height=2.05cm,inner sep=5pt,font=\small},
  arrow/.style={-{Stealth[length=2mm]},semithick},
  node distance=7mm]
\node[box] (anchor) {\textbf{Real anchor}\\[4pt]
  $S\subseteq F$\\$|S|=h$};
\node[box,right=of anchor] (reference) {\textbf{Sampled reference}\\[4pt]
  $h-|Q|$ real facilities\\
  $a$ leaves and one\\permanent dummy};
\node[box,right=of reference] (output) {\textbf{Real completion}\\[4pt]
  remove $r=m-|Q|$ more\\
  add at most $a+c=m$\\
  keep at most $h$ facilities};
\draw[arrow] (anchor)--(reference);
\draw[arrow] (reference)--(output);
\node[below=8mm of reference,align=center,text width=14.8cm,font=\small]
  {$\displaystyle h-|Q|-r+a+c=h$\quad before merging repeated choices.\\[3pt]
   Ball destinations: at most $a$; destinations without a ball: at most $c=m-a$.};
\end{tikzpicture}
\caption{The facility count in completion.  The reference contains virtual
leaves, but the output contains only members of $F$.  A chosen facility
in a ball replaces its leaf; a nonempty real solution dominates the
permanent dummy.  The removals are distinct, while additions are set
unions, so overlaps and repeated destinations only reduce the count.
The figure describes accounting, not a Euclidean embedding or a
disjointness assumption on facility balls.}
\label{fig:completion-budget}
\end{figure}
